\section{Introduction}

Stable communication is essential during disasters and large-scale power outages, when affected populations and emergency responders depend on cellular networks to obtain and exchange timely information. When grid power is unavailable, radio access networks must rely on finite backup-energy resources to maintain service. In an O-RAN deployment, Radio Unit (RU) operation is therefore critical: an RU that depletes its battery can no longer provide radio access to users within its coverage area~\cite{oranScArchitecture}.

This scenario implies a fundamental trade-off between energy consumption and service availability. Keeping all RUs active maximizes coverage but accelerates battery depletion. In contrast, placing selected RUs into a low-power sleep state conserves energy and extends battery availability, though it may reduce the proportion of users served. Sleep-mode mechanisms are established methods for reducing cellular-network energy consumption \cite{wu2015sleep}; however, during an outage, energy savings alone are insufficient. A policy that preserves battery capacity but does not sustain an acceptable service level cannot be considered dependable for the modeled scenario.

In this work, we study whether simple, predefined RU sleep policies can balance emergency connectivity and battery lifetime during a power outage. Three baseline policies are evaluated: Always Active, in which all RUs remain operational throughout the outage; Staggered Active, where RUs alternate periods of activity; and Threshold Staggered Active, where a predefined threshold determines when RUs rotate their activity. These policies serve as reference points for assessing the trade-offs and limitations of straightforward sleep scheduling. The central research questions are: \emph{How do predefined RU sleep policies influence emergency connectivity and RU battery lifetime in an O-RAN network during a power outage?} \emph{Do these policies meet the configured emergency-service target, or are more adaptive outage-management strategies required?} This work is a controlled, policy-level simulation study rather than a proposal for an optimal scheduling algorithm, a physical-layer-accurate analysis, or a validation of a specific hardware deployment.

This work makes four main contributions:
\begin{itemize}
    \item We formulate a simplified, scenario-specific model of an O-RAN power outage that captures RU battery depletion, active and sleep states, stochastic link quality, capacity-constrained user association, and emergency-service measures.
    \item We conducted controlled simulations of three predefined RU operating policies: Always Active, Staggered Active, and Threshold Staggered Active. The evaluation considers emergency QoS, RU battery-depletion time, and SLA-based network lifetime as complementary measures.
    \item Within the stated model configuration, results show that extending RU battery lifetime does not always ensure reliable emergency service. Periodic staggering maximizes battery extension but often fails to meet the service target early. In contrast, threshold-triggered staggering maintains the service lifetime of the Always Active policy, but offers only a modest battery benefit.
    \item To perform our evaluations, we developed a modular, discrete-time simulation system. This simulator separates policy decisions from the network environment, execution procedure, and metric collection, enabling reproducible comparison of alternative RU control policies.
\end{itemize}

The results show that the battery benefits of predefined sleep policies do not necessarily ensure stable emergency operation. Periodic staggered sleeping achieves the greatest extension of RU battery lifetime but often fails to maintain the configured service target from the start of the outage. Threshold-triggered staggering maintains an SLA-based network lifetime similar to the always-active baseline, yet causes only a modest improvement in battery lifetime. Therefore, under the modeled conditions, the evaluated fixed policies serve as useful reference points but do not provide a sufficiently sound foundation for emergency outage management. Their limitations highlight the need for more adaptive policy designs. This conclusion is specific to the simplified model and experimental configuration employed in this study and does not establish the suitability of these policies for all real-world deployments.

The remainder of this paper is organized as follows. Section~II presents background and related work. Section~III defines the system model and evaluated policies. Section~IV describes the simulation methodology and experimental configuration. Section~V discusses the simulation results, and Section~VI concludes the paper.
